\subsubsection*{Boolean Functions and Representation}

% BEGIN GENERATED ARTIFACT: def:boolean-function-representation-propositional-logic
\begin{definitionbox}{Definition (Representation of a Boolean Function)}
\begin{definition}[Representation of a Boolean Function]
\label{def:boolean-function-representation-propositional-logic}
Let \(f:\mathbb{B}^n\to\mathbb{B}\) be a Boolean function. A formula
\(\varphi\in\WFF\) with variables among \(P_1,\ldots,P_n\) \emph{represents}
\(f\) if for every truth assignment \(v\),
\[
\widehat v(\varphi)=f(v(P_1),\ldots,v(P_n)).
\]
\end{definition}
\end{definitionbox}

\begin{remark*}[Standard quantified statement]
\[
\operatorname{Represents}(\varphi,f;P_1,\ldots,P_n)
\Longleftrightarrow
(\forall v:\Prop\to\mathbb{B})
\bigl(\widehat v(\varphi)=f(v(P_1),\ldots,v(P_n))\bigr).
\]
\end{remark*}

\begin{remark*}[Definition predicate reading]
\(\operatorname{Represents}(\varphi,f;P_1,\ldots,P_n)\) means that
\(\varphi\) has exactly the same truth table as the Boolean function \(f\) on
the listed variables.
\end{remark*}

\begin{remark*}[Interpretation]
A formula represents a Boolean function when the formula computes the same
truth value as the function for every input row. This turns a semantic table of
values into a syntactic object inside the propositional language.
\end{remark*}

\begin{dependencies}
\begin{itemize}
  \item \hyperref[def:boolean-function-propositional-logic]{Boolean Function}
  \item \hyperref[def:truth-function-induced-by-formula]{Truth Function Induced by a Formula}
  \item \hyperref[def:truth-assignment-propositional-logic]{Truth Assignment}
\end{itemize}
\end{dependencies}
% END GENERATED ARTIFACT: def:boolean-function-representation-propositional-logic

% BEGIN GENERATED ARTIFACT: thm:dnf-representation-of-boolean-functions-propositional-logic
\begin{theorembox}{Theorem (DNF Representation of Boolean Functions)}
\begin{theorem}[DNF Representation of Boolean Functions]
\label{thm:dnf-representation-of-boolean-functions-propositional-logic}
Every Boolean function \(f:\mathbb{B}^n\to\mathbb{B}\) is represented by a
formula in disjunctive normal form using variables \(P_1,\ldots,P_n\).

\noindent\hyperref[prf:dnf-representation-of-boolean-functions-propositional-logic]{Go to proof.}
\end{theorem}
\end{theorembox}

\begin{remark*}[Standard quantified statement]
\[
(\forall f:\mathbb{B}^n\to\mathbb{B})(\exists\varphi\in\WFF)
\bigl(\varphi\text{ is in DNF}\land
\operatorname{Represents}(\varphi,f;P_1,\ldots,P_n)\bigr).
\]
\end{remark*}

\begin{remark*}[Predicate reading]
Every \(n\)-ary Boolean function has a propositional formula representative, and
one may choose the representative to be in disjunctive normal form.
\end{remark*}

\begin{remark*}[Interpretation]
A truth table specifies which input rows make the Boolean function true. For
each true row, form the corresponding conjunction of literals; the disjunction
of those row conjunctions represents the entire function.
\end{remark*}

\begin{dependencies}
\begin{itemize}
  \item \hyperref[def:boolean-function-representation-propositional-logic]{Representation of a Boolean Function}
  \item \hyperref[def:canonical-disjunctive-normal-form-propositional-logic]{Canonical Disjunctive Normal Form}
  \item \hyperref[thm:canonical-disjunctive-normal-form-propositional-logic]{Canonical Disjunctive Normal Form Theorem}
  \item \hyperref[def:disjunctive-normal-form-propositional-logic]{Disjunctive Normal Form}
\end{itemize}
\end{dependencies}
% END GENERATED ARTIFACT: thm:dnf-representation-of-boolean-functions-propositional-logic
